\begingroup
\scriptsize
\setlength{\tabcolsep}{2pt}
\setlength{\LTleft}{0pt}
\setlength{\LTright}{0pt}
\begin{longtable}{>{\raggedright\arraybackslash}p{0.14\textwidth}>{\raggedright\arraybackslash}p{0.16\textwidth}>{\centering\arraybackslash}p{0.06\textwidth}>{\centering\arraybackslash}p{0.12\textwidth}>{\raggedright\arraybackslash}p{0.43\textwidth}}
\caption{Complete ADEMP mechanism specification.}\label{tab:supp_ademp_design}\\
\toprule
Group & Mechanism & Stays & Index/later counts & Defining feature\\
\midrule
\endfirsthead
\multicolumn{5}{l}{\tablename\ \thetable\ continued}\\
\toprule
Group & Mechanism & Stays & Index/later counts & Defining feature\\
\midrule
\endhead
\midrule\multicolumn{5}{r}{Continued on next page}\\
\endfoot
\bottomrule
\endlastfoot
Sample size and density & Large dense & 600 & (8, 16) / (8, 16) & Correct common component, independent errors, persistent level, and dense non-informative observation.\\
Sample size and density & Small dense & 240 & (8, 16) / (8, 16) & The ideal mechanism with 240 stays to isolate sample-size sensitivity.\\
Sample size and density & Large sparse & 600 & (2, 5) / (4, 8) & The ideal mechanism with only two to five index observations to isolate observation sparsity.\\
Observation process & Serial dependence & 600 & (6, 14) / (6, 14) & A persistent level is observed with a Gaussian-copula AR(1) residual process on a 15-minute latent grid.\\
Observation process & Integer rounding & 600 & (6, 14) / (6, 14) & Observed and oracle MAP quantiles are rounded to integer mmHg, creating ties that require interval calibration.\\
Observation process & Informative monitoring & 600 & (4, 22) / (4, 22) & Observation probability increases for lower latent levels and after a recently observed low MAP.\\
Observation process & Cluster size association & 600 & (3, 20) / (3, 20) & Cluster size is associated with the latent stay level even though observation times are otherwise non-informative.\\
Trajectory structure & Omitted common time & 600 & (6, 14) / (6, 14) & The true common quantile contains nonlinear circadian and quadratic time structure omitted from the fitted common component.\\
Trajectory structure & Persistent level with shape & 600 & (8, 16) / (8, 16) & Persistent random intercepts coexist with stay-specific linear time shape, favoring a correctly targeted level-plus-slope update.\\
Trajectory structure & Treatment feedback & 600 & (6, 14) / (6, 14) & An early low lower tail triggers a later MAP-raising intervention whose effect decays over time.\\
Signal controls & Transient displacement & 600 & (6, 14) / (6, 14) & A large index-only shift disappears later, providing a negative control under serial residual dependence.\\
Signal controls & Serial null & 600 & (6, 14) / (6, 14) & No persistent level, index shift, or random shape is present; serial residual dependence tests whether tuning suppresses noise-only updates.\\
Signal controls & Weak persistent level & 600 & (8, 16) / (8, 16) & A small but persistent stay-specific lower-tail level tests power and shrinkage under a weak signal.\\
Signal controls & Heavy tailed residual & 600 & (8, 16) / (8, 16) & A persistent level is observed with independent Student-t3 marginal residuals centered at their exact 0.10 quantile.\\
\end{longtable}
\endgroup
